\begin{proof}

\begin{enumerate}[label={$\langle1\rangle$\arabic*.},leftmargin=*]
\item We derive the needed holomorphic invariants from the local models.
The commented-out holomorphic-invariants notes in the supplied
Liu--Rollenske source are not used as established results.

\item In the Weierstrass equation of \cite[Proposition labelled
\texttt{prop: surface facts}]{LiuRollenske2026}, make the explicit
coordinate comparison
\[
 t_{\rm new}=\frac{27}{4}(t_{\rm old}^{-1}-1),\quad
 x_{\rm new}=x_{\rm old}/t_{\rm old}^2,\quad
 y_{\rm new}=y_{\rm old}/t_{\rm old}^3.
\]
It gives $y^2=x^3+tx+t$, with section $P=(-1,i)$; the multiple
value is now infinity, the cusp is at zero and the remaining nodal
value is $-27/4$. Consider the explicit deformation
\begin{equation}\hypertarget{tag-01A8}{}\label{cdp:eq:E14-elliptic-family}
 y^2=x^3+tx+t+\epsilon,\qquad
 P_\epsilon=(-1,b),\quad b^2=\epsilon-1,
 \quad b(0)=i.
\end{equation}
Its finite discriminant is $4t^3+27(t+\epsilon)^2$.
The nodal value near $-27/4$ remains simple and holomorphically
labelled. The other two nodes coalesce to the type-II cusp at
$\epsilon=0$. Infinity keeps type $III^*$; on its order-four
good-reduction cover the equation is
\[
 y'^2=x'^3+x'+r^2+\epsilon r^6,
 \qquad P'_{\epsilon}=(-r^2,br^3).
\]

\item These equations give a simultaneous construction of the compact
threefolds, including the central one. On a disc containing the cusp
and its two nearby nodal values, the total elliptic surface in
\eqref{cdp:eq:E14-elliptic-family} is smooth over $\epsilon$:
$\partial(y^2-x^3-tx-t-\epsilon)/\partial t=-(x+1)$ is a unit
near the cusp. Translation by $-2P_\epsilon$ is an automorphism
there. The line bundle $\mathcal O(P_\epsilon-O)$ and its lifted
translation vary holomorphically; its multiplier has a zero only
at the separate simple nodal value chosen for the hexagonal filling.
For a uniformly small nonzero scale its contraction on the
complement of the zero section is uniform on the compact elliptic
fibres. The quotient is consequently free, properly discontinuous,
proper over the disc, and smooth over the parameter. This quotient
itself fills the cusp and the nearby rank-one nodal fibres. At the
separate node the ordinary rank-two Mumford charts vary
holomorphically. At infinity the fixed order-four torsion translation
and good-reduction comparison vary holomorphically in the displayed
equation. Gluing along their marked punctured collars therefore
constructs a proper holomorphic family over $(\epsilon,u)$ whose
central scale slice is $Y(u)$. Every member has the central
differentiable type by Ehresmann. The supplied topology theorem
\cite[Theorem labelled \texttt{thm: main}]{LiuRollenske2026} thus
identifies these differentiable manifolds with $S^6$.

\item Write $\mathfrak B=\mathbb P(1,4)$ and retain its convention
$\rho^*\mathcal O(1)=\mathcal O(4)$. On the central threefold the
canonical form on the contracting quotient is the elliptic-surface
two-form wedged with $dv/v$. At the cusp it extends without a zero
or pole: from $t=(y^2-x^3)/(1+x)$ one gets
\[
 dt\wedge\frac{dx}{2y}=-\frac{dx\wedge dy}{1+x}.
\]
The finite-cover and hexagonal calculations are the same explicit
coordinate calculations as in \ref{cdp:st:E13}. Since the elliptic surface has
$K_S=\pi^*\mathcal O_{\mathbb P^1}(-1)$, these comparisons give
\begin{equation}\hypertarget{tag-01A7}{}\label{cdp:eq:E14-canonical}
 K_Y=f^*\mathcal O_{\mathbb P^1}(-1)
     =\bar f^*\mathcal O_{\mathfrak B}(-4).
\end{equation}
The calculation is simultaneous in $\epsilon$. Adjunction on every
full fibre gives trivial dualizing sheaf. \CompendiumRef{II}{cdp:st:E4} gives $h^0=1$;
duality gives $h^2=1$; constancy of its Euler characteristic in the
proper flat family gives $h^1=2$. Thus, also on the two-dimensional
base $(t,\epsilon)$, Grauert's theorem gives locally free $R^qf_*\mathcal O$
with base change, for $q=0,1,2$.

\item The multiplicative connection
\[
 \alpha=\frac{y+b}{x+1}\frac{dx}{2y},\qquad
 \nabla(hs)=(dh+\alpha h)s
\]
has precisely the residues of the rational section $s$ with divisor
$P_\epsilon-O$. The same proof and good-reduction substitution as
in \ref{cdp:st:E13} give the logarithmic Lie bundle
$E=\mathcal O\oplus\mathcal O(-1)$, including at the central
cusp. There the vector dual to $dx/(2y)$ is the derivation
$2y\partial_x+3x^2\partial_y$ on the cuspidal fibre, and its
horizontal lift is well defined because
$(y+b)/(x+1)$ is regular near the singular locus.

\item We justify the identification with $G=R^1\bar f_*\mathcal O$
at the cusp without assuming a limiting-period rule there.
The inverse pairing associated with the invariant alternating
class gives $G\simeq E$ over the two-dimensional $(t,\epsilon)$
base away from its single cuspidal point: all remaining points
are smooth, nodal or the finite-cover points already computed in
\ref{cdp:st:E13}. Both $G$ and $E$ are locally free across that point, by the
preceding cohomological-flatness and connection calculations.
The matrix entries of the isomorphism and its inverse extend
across that codimension-two set by Hartogs' theorem. Their products
are the identity on a dense open set, hence everywhere. Thus
\[
 G=\mathcal O\oplus\mathcal O(-1),\qquad
 R^2\bar f_*\mathcal O=\mathcal O(-1)
\]
on the central root curve as well; the second equality also follows
directly from relative duality and \eqref{cdp:eq:E14-canonical}.

\item Let $D=C+W$ be the sum of the cuspidal and hexagonal fibres.
For the local tangent calculation choose coordinates in which
$t=y^2-x^3$. They are obtained from the earlier formula by the
invertible change $\widetilde x=x/(1+x)^{1/3}$,
$\widetilde y=y/(1+x)^{1/2}$, with the branches equal to one at
zero, and then renaming the two new coordinates. The logarithmic
tangent module is freely generated by
\[
 2y\partial_x+3x^2\partial_y,\quad
 2x\partial_x+3y\partial_y,\quad\partial_z,
\]
where $z$ is a coordinate along the elliptic singular curve. The
second generator maps to $6t\partial_t$;
the first and third generate the relative logarithmic kernel.
Consequently the global exact sequence is
\[
 0\longrightarrow\bar f^*E\longrightarrow T_Y(-\log D)
 \longrightarrow\bar f^*\mathcal O(-3)\longrightarrow0,
\]
because $T_{\mathfrak B}(-C-W)=\mathcal O(5-8)$.
The vertical cohomology is $(1,1,0,0)$. Every line in the
horizontal Leray terms is $\mathcal O(-3)$ or $\mathcal O(-4)$,
which has no cohomology on $\mathbb P(1,4)$. Hence
\begin{equation}\hypertarget{tag-01AA}{}\label{cdp:eq:E14-log-coho}
 h^\bullet(T_Y(-\log D))=(1,1,0,0).
\end{equation}

\item The quotient at $W$ is $\eta_*K_{dP_6}$ with cohomology $(0,0,1)$.
At $C$, the map $v\mapsto v(t)|_C$ identifies the quotient with
the Jacobian ideal
\[
 J=(x^2,y)\subset\mathcal O_C.
\]
Indeed its kernel is precisely the logarithmic tangent module and
$N_{C/Y}=\mathcal O_C$. Put $Z=V(J)\subset C$. We now compute
the actual restriction map
\begin{equation}\hypertarget{tag-01A9}{}\label{cdp:eq:E14-jet-map}
 H^1(C,\mathcal O_C)\longrightarrow H^1(Z,\mathcal O_Z).
\end{equation}

\item In the deformation \eqref{cdp:eq:E14-elliptic-family} the critical
base near the cusp is the smooth curve
\[
 \Gamma:\quad t=-3x^2,\qquad
 \epsilon=x^2(3+2x).
\]
Its map to the $\epsilon$-disc is finite flat of degree two; its
central fibre is $\operatorname{Spec}\mathbb C[x]/(x^2)$.
The lifted translation fixes its $x$-coordinate. This follows on
the punctured curve because translation fixes the unique node of
each nodal cubic, and then follows at zero by holomorphic
continuation. On the line coordinate it is
$(x,v)\mapsto(x,c(x)v)$ with $c$ a holomorphic nowhere-zero
contraction. Thus its quotient critical locus is exactly the
elliptic family
\[
 \mathcal Z=(\Gamma\times\mathbb C^*)/
       \langle(x,v)\mapsto(x,c(x)v)\rangle.
\]
In particular $H^0(Z,\mathcal O_Z)$ and
$H^1(Z,\mathcal O_Z)$ are both free rank-one modules over
$A=\mathbb C[x]/(x^2)$, and $H^2(Z,\mathcal O_Z)=0$.

\item Restriction in this proper family is a holomorphic bundle map
$G|_\Gamma\to R^1(\mathcal Z/\Gamma)_*\mathcal O_{\mathcal Z}$.
On $\Gamma\setminus\{0\}$ the calculation of \ref{cdp:st:E13} identifies it,
after a nonzero scalar change of target frame, with the row
\[
 \left(1,-\frac{b}{1+x}\right).
\]
The same frame can be chosen at zero. To check this last point,
normalize $C$. Its normalization is a $\mathbb P^1$-bundle over
the elliptic singular curve $E_0$, and its normalization quotient
is $\mathcal O_{E_0}$: the first jet at the cusp transforms with
derivative one under translation on the additive smooth locus.
The normalization exact sequence, together with
$h^1(\mathcal O_C)=2$, $h^2(\mathcal O_C)=1$ and
$h^\bullet(\mathcal O_{\widetilde C})=(1,1,0)$, implies
that $H^1(\mathcal O_C)\to H^1(\mathcal O_{\widetilde C})$
is onto. Restriction from the ruled normalization to its singular
section identifies the latter group with $H^1(\mathcal O_{E_0})$.
Thus restriction to $E_0$ is onto. The common scalar in the row
above cannot have a zero at zero, since both entries would then
vanish and this restriction would be zero. It is a unit. This
proves the asserted holomorphic target-frame comparison.

\item The equality of these holomorphic maps on the punctured curve
extends over all of $\Gamma$, and in particular over its
nonreduced central fibre. Hence \eqref{cdp:eq:E14-jet-map} is,
in the specified bases, the complex-linear map
\[
 \mathbb C^2\longrightarrow A,\qquad
 (a,d)\longmapsto a+d(-i+ix).
\]
Its matrix in $1,x$ has determinant $i$, so it is an isomorphism.
This is the required jet comparison, not an inference from the
rank at nearby reduced fibres.

\item The exact sequence $0\to J\to\mathcal O_C\to\mathcal O_Z\to0$
now gives
\[
 h^\bullet(C,J)=(0,1,1).
\]
Together with \eqref{cdp:eq:E14-log-coho} and the hexagonal quotient,
the tangent exact sequence gives
$h^\bullet(T_Y)=(1,2,2,0)$.

\item It remains to verify that the actual two-parameter family has
bijective Kodaira--Spencer map. In the above local trivialization
of its total space the map to the base is
\[
 f_{\epsilon}=\frac{y^2-x^3-\epsilon}{1+x}.
\]
The derivative $\partial_\epsilon f$ has image $-1+x$ in
$\mathcal O_C/J=\mathbb C[x]/(x^2)$ along the singular curve.
If the ambient first-order deformation were trivial, every
first-order variation of its map would be a global section of
$f^*T_{\mathbb P^1}$, hence the pullback of a projective vector
field, because $f_*\mathcal O_Y=\mathcal O_{\mathbb P^1}$.
Its local image in this Jacobian quotient is constant. Changing
the local ambient trivialization adds $df(v)$ and hence adds
zero in that quotient. The nonzero coefficient of $x$ therefore
proves that the $\epsilon$-direction is nontrivial and cannot be
cancelled by the scale direction, which does not change this local
map. With $\epsilon=0$, a normalized singular elliptic curve is
$\mathbb C^*/\langle c_0u\rangle$; its modulus has nonzero
$u$-derivative for sufficiently small $u$. A trivial deformation
of the total space preserves its anticanonical pencil from
\eqref{cdp:eq:E14-canonical}, its three distinguished special values,
and that elliptic singular curve. Hence the scale direction is
nontrivial too. The two directions span $H^1(T_Y)$.
The inverse-function argument of \ref{cdp:st:E13} consequently makes the
Kuranishi ideal zero and proves miniversality. \ref{cdp:st:E6} and \ref{cdp:st:E1} give the
remaining assertions.

\end{enumerate}
\end{proof}
